Las llaves M1 y M2 de la \textit{buck} se implementan con dos MOSFET de forma tal que el circuito conmute. Cuando M1 está cerrada y M2 está abierta, la corriente circula por la malla general y el inductor tiene una polaridad positiva en el nodo de conmutación. Pero, cuando M1 se abre, M2 permanece momentáneamente abierta, abriendo el circuito por completo. Para ello, se coloca un diodo de rueda libre en paralelo al MOS que le otorga un camino alternativo a la corriente para que circule por la malla y permitiendo que el inductor invierta su polaridad (conmuta).

En estos casos, en el momento de la conmutación se producen corrientes que circulan hacia la fuente mediante el MOS M1 que pudieran dañar el diodo intrínseco del mismo, por lo que se le conecta en paralelo otro diodo de rueda libre para evitarlo.

Se escogen diodos Schottky para esta función debido a su baja caída de tensión directa, lo que reduce las pérdidas de conducción durante los intervalos en los que circula la corriente de rueda libre. Además, presentan una recuperación inversa prácticamente nula, característica importante en una fuente conmutada, ya que disminuye los picos de corriente y las pérdidas asociadas al cambio de estado. De esta manera, se mejora la eficiencia del convertidor y se reducen las exigencias eléctricas sobre los MOSFET durante la conmutación.